% GENERATED FILE. Edit canonical lessons, then run npm run generate:books.
% canonical-source-hash: fnv1a64:9c5103f07b2c2f18
% canonical-lessons: JA-C158-kagi, JA-C158-kesu, JA-C158-denki, JA-C158-akeru, JA-C158-kurai

\chapter{Keys, Windows and Lights}
\label{ch:ja-keys-windows-and-lights}
\hlchaptermodality{158}

\begin{chapteropening}
\textbf{By the end of this chapter:} \emph{I can talk about locking up, opening windows and switching off lights.}

Complete the last lesson of chapter 158: I can talk about locking up, opening windows and switching off lights.
\end{chapteropening}

\section[\texorpdfstring{\ja{かぎ} (kagi)}{かぎ (kagi)}]{\ja{かぎ} (kagi) — a key; a lock}

\label{lesson:JA-C158-kagi}



\begin{quote}
Before the new one: say the Japanese for to hang, then the Japanese for to put.
\end{quote}

\subsection*{You'll want to know: \ja{かぎ}}
\textbf{\ja{かぎ}} — \emph{kagi} — \textquotedblleft{}a key; a lock\textquotedblright{}.

\begin{quote}
\textbf{\ja{かぎを} \ja{かばんに} \ja{いれます。}} — \emph{kagi o kaban ni iremasu} — \textquotedblleft{}I put the key in my bag.\textquotedblright{}
\end{quote}

\begin{cousinweb}
\textbf{\ja{かぎ}} (\emph{kagi}) is a native Japanese word (wago), inherited rather than borrowed from Chinese; its kanji was attached later as a kun'yomi, and its older origin is not settled (JMdict, KANJIDIC2).
\end{cousinweb}

\begin{grammarlens}[title={What you've built}]
One more word for everyday life at A2.
\end{grammarlens}

\subsection*{Your turn}
\begin{itemize}
\raggedright
  \item \emph{Say it:} \emph{kagi}
  \item \emph{Say it:} \emph{kagi}, once more
  \item \emph{Say it:} \emph{kagi o kaban ni iremasu} — I put the key in my bag.
  \item \emph{Recall:} say the Japanese for to hang, then the Japanese for to put, then say \emph{kagi} again
\end{itemize}

\begin{tcolorbox}[breakable,colback=teal!4,colframe=teal!35!black,title={Before you move on}]
What does \ja{かぎ} mean? (\textquotedblleft{}A key; a lock\textquotedblright{}.) Say it once more.
\end{tcolorbox}
\section[\texorpdfstring{\ja{けす} (kesu)}{けす (kesu)}]{\ja{けす} (kesu) — to switch off; to put out}

\label{lesson:JA-C158-kesu}



\begin{quote}
Before the new one: say the Japanese for to put, then the Japanese for a key.
\end{quote}

\subsection*{You'll want to know: \ja{けす}}
\textbf{\ja{けす}} — \emph{kesu} — \textquotedblleft{}to switch off; to put out\textquotedblright{}.

\begin{quote}
\textbf{\ja{あかりを} \ja{けします。}} — \emph{akari o keshimasu} — \textquotedblleft{}I switch off the lamp.\textquotedblright{}
\end{quote}

Polite \textbf{\ja{けします}}, \emph{keshimasu}; a godan verb.

\begin{cousinweb}
\textbf{\ja{けす}} (\emph{kesu}) is a native Japanese word (wago), inherited rather than borrowed from Chinese; its kanji was attached later as a kun'yomi, and its older origin is not settled (JMdict, KANJIDIC2).
\end{cousinweb}

\begin{grammarlens}[title={What you've built}]
One more word for everyday life at A2.
\end{grammarlens}

\subsection*{Your turn}
\begin{itemize}
\raggedright
  \item \emph{Say it:} \emph{kesu}
  \item \emph{Say it:} \emph{kesu}, once more
  \item \emph{Say it:} \emph{akari o keshimasu} — I switch off the lamp.
  \item \emph{Recall:} say the Japanese for to put, then the Japanese for a key, then say \emph{kesu} again
\end{itemize}

\begin{tcolorbox}[breakable,colback=teal!4,colframe=teal!35!black,title={Before you move on}]
What does \ja{けす} mean? (\textquotedblleft{}To switch off; to put out\textquotedblright{}.) Say it once more.
\end{tcolorbox}
\section[\texorpdfstring{\ja{でんき} (denki)}{でんき (denki)}]{\ja{でんき} (denki) — electric light; electricity}

\label{lesson:JA-C158-denki}



\begin{quote}
Before the new one: say the Japanese for a key, then the Japanese for to switch off.
\end{quote}

\subsection*{You'll want to know: \ja{でんき}}
\textbf{\ja{でんき}} — \emph{denki} — \textquotedblleft{}electric light; electricity\textquotedblright{}.

\begin{quote}
\textbf{\ja{へやの} \ja{でんきを} \ja{けします。}} — \emph{heya no denki o keshimasu} — \textquotedblleft{}I switch off the light in the room.\textquotedblright{}
\end{quote}

\begin{cousinweb}
\textbf{\ja{でんき}} (\emph{denki}) is Sino-Japanese (kango): two Chinese characters read with their on'yomi, readings borrowed from Chinese. They are \emph{den}, \textquotedblleft{}electricity\textquotedblright{} (Middle Chinese \emph{denH}) and \emph{ki}, \textquotedblleft{}air, spirit\textquotedblright{} (\emph{khj+jH}) (KANJIDIC2; Middle Chinese in Baxter's transcription).
\end{cousinweb}

\begin{grammarlens}[title={What you've built}]
One more word for everyday life at A2.
\end{grammarlens}

\subsection*{Your turn}
\begin{itemize}
\raggedright
  \item \emph{Say it:} \emph{denki}
  \item \emph{Say it:} \emph{denki}, once more
  \item \emph{Say it:} \emph{heya no denki o keshimasu} — I switch off the light in the room.
  \item \emph{Recall:} say the Japanese for a key, then the Japanese for to switch off, then say \emph{denki} again
\end{itemize}

\begin{tcolorbox}[breakable,colback=teal!4,colframe=teal!35!black,title={Before you move on}]
What does \ja{でんき} mean? (\textquotedblleft{}Electric light; electricity\textquotedblright{}.) Say it once more.
\end{tcolorbox}
\section[\texorpdfstring{\ja{あける} (akeru)}{あける (akeru)}]{\ja{あける} (akeru) — to open up; to open (a door, a window)}

\label{lesson:JA-C158-akeru}



\begin{quote}
Before the new one: say the Japanese for to switch off, then the Japanese for electric light.
\end{quote}

\subsection*{You'll want to know: \ja{あける}}
\textbf{\ja{あける}} — \emph{akeru} — \textquotedblleft{}to open up; to open (a door, a window)\textquotedblright{}.

\begin{quote}
\textbf{\ja{あさ} \ja{まどを} \ja{あけます。}} — \emph{asa mado o akemasu} — \textquotedblleft{}I open the window in the morning.\textquotedblright{}
\end{quote}

Polite \textbf{\ja{あけます}}, \emph{akemasu}; an ichidan verb.

\begin{cousinweb}
\textbf{\ja{あける}} (\emph{akeru}) is a native Japanese word (wago), inherited rather than borrowed from Chinese; its kanji was attached later as a kun'yomi, and its older origin is not settled (JMdict, KANJIDIC2).
\end{cousinweb}

\begin{grammarlens}[title={What you've built}]
One more word for everyday life at A2.
\end{grammarlens}

\subsection*{Your turn}
\begin{itemize}
\raggedright
  \item \emph{Say it:} \emph{akeru}
  \item \emph{Say it:} \emph{akeru}, once more
  \item \emph{Say it:} \emph{asa mado o akemasu} — I open the window in the morning.
  \item \emph{Recall:} say the Japanese for to switch off, then the Japanese for electric light, then say \emph{akeru} again
\end{itemize}

\begin{tcolorbox}[breakable,colback=teal!4,colframe=teal!35!black,title={Before you move on}]
What does \ja{あける} mean? (\textquotedblleft{}To open up; to open (a door, a window)\textquotedblright{}.) Say it once more.
\end{tcolorbox}
\section[\texorpdfstring{\ja{くらい} (kurai)}{くらい (kurai)}]{\ja{くらい} (kurai) — dark}

\label{lesson:JA-C158-kurai}



\begin{quote}
Before the new one: say the Japanese for electric light, then the Japanese for to open up.
\end{quote}

\subsection*{You'll want to know: \ja{くらい}}
\textbf{\ja{くらい}} — \emph{kurai} — \textquotedblleft{}dark\textquotedblright{}.

\begin{quote}
\textbf{\ja{この} \ja{へやは} \ja{くらいです。}} — \emph{kono heya wa kurai desu} — \textquotedblleft{}This room is dark.\textquotedblright{}
\end{quote}

\begin{cousinweb}
\textbf{\ja{くらい}} (\emph{kurai}) is a native Japanese word (wago), inherited rather than borrowed from Chinese; its kanji was attached later as a kun'yomi, and its older origin is not settled (JMdict, KANJIDIC2).
\end{cousinweb}

\begin{grammarlens}[title={What you've built}]
One more word for everyday life at A2.
\end{grammarlens}

\subsection*{Your turn}
\begin{itemize}
\raggedright
  \item \emph{Say it:} \emph{kurai}
  \item \emph{Say it:} \emph{kurai}, once more
  \item \emph{Say it:} \emph{kono heya wa kurai desu} — This room is dark.
  \item \emph{Recall:} say the Japanese for electric light, then the Japanese for to open up, then say \emph{kurai} again
\end{itemize}

\begin{tcolorbox}[breakable,colback=teal!4,colframe=teal!35!black,title={Before you move on}]
What does \ja{くらい} mean? (\textquotedblleft{}Dark\textquotedblright{}.) Say it once more.
\end{tcolorbox}
